\documentclass[10pt,a4paper]{amsart}
\usepackage[margin=19mm]{geometry}
\usepackage{amsmath,amssymb,mathtools}
\usepackage[scr=rsfs,cal=euler]{mathalfa}
\usepackage{xfrac}
\usepackage[unicode,hidelinks,bookmarks=false]{hyperref}
\newcommand{\ie}{i.e.\ }
\newcommand{\cf}{cf.\ }
\newcommand{\resp}{respectively\ }
\newcommand{\Subsection}{\subsection}
\providecommand{\nobreakdash}{\nobreak\hbox{-}}
\setlength{\emergencystretch}{2em}
\multlinegap0pt
\usepackage[shortlabels]{enumitem}
\usepackage{booktabs}
\usepackage{nicefrac}
\usepackage{xcolor}
\usepackage{multirow}
\usepackage{caption}
\usepackage{bbm}
\usepackage{breqn}
\usepackage{bm}
\usepackage{algorithm}\usepackage{algpseudocode}
\newtheorem{lemma}{Lemma}
\newcommand{\E}{\mathbb{E}}
\DeclareMathOperator{\Tr}{Tr}
\newcommand{\bas}[1]{\begin{align*}#1\end{align*}}
\newcommand{\bb}[1]{\left(#1\right)}
\newcommand{\bbb}[1]{\left[#1\right]}
\newcommand{\be}{\pmb{e}}
\newcommand{\bk}{\color{black}}
\newcommand{\bt}{\pmb{t}}
\newcommand{\bu}{\pmb{u}}
\newcommand{\bx}{\pmb{x}}
\newcommand{\norm}[1]{\left\|#1\right\|}
\newcommand{\nt}{\|\bt\|}
\newcommand{\nts}{\|\bt\|^2}
\newcommand{\rd}{\color{red}}
\title{Nonparametric Evaluation of Noisy ICA Solutions}
\author{Syamantak Kumar, Purnamrita Sarkar, Peter Bickel,, Derek Bean}
\date{Source: arXiv:2401.08468v4 (CC BY 4.0)}
\begin{document}
\maketitle

\noindent\emph{Source and licence.} Nonparametric Evaluation of Noisy ICA Solutions. Version arXiv:2401.08468v4 (CC BY 4.0), distributed by arXiv under the Creative Commons Attribution 4.0 International licence, http://creativecommons.org/licenses/by/4.0/, as recorded in the arXiv licence metadata for this exact version and shown on the abstract page https://arxiv.org/abs/2401.08468v4. Full licence notice: \texttt{licenses/arxiv-2401.08468-CC-BY-4.0.txt}.\par\smallskip

\noindent\textit{Editorial scope.} Counted unit: the article's lemma lem:scorerhs (source0.tex lines 1674--1679). The article's complete original proof of it is delivered with it (source0.tex lines 1680--1722, one \texttt{proof} environment of 43 lines). No mathematics is written, repaired or re-derived: the statement and the proof are the article's own lines, in the article's own notation, and no retained line is substituted or re-flowed.  Retained uncounted as the source-local context the proof needs: lines 1588--1595.  Nothing was omitted from any retained span, so the package carries no omission record.  No retained line contains a citation, so the wrapper carries no bibliography and no bibliography entry.  No retained line contains a figure environment, an image-inclusion command or drawing code, so no image is delivered and the package declares no asset dependency.  Editorial formatting only, in addition: an AMS \texttt{amsart} preamble that restates the article's own declarations and macros used by the retained lines, namely the environments \texttt{lemma} on separate counters, together with the standard packages those lines need.  The retained environments print in this document as Lemma 1 in order of appearance, whereas the article prints the same items with the numbering of its own full document.  The complete native source of the article as distributed in the arXiv e-print for this exact version is bundled unchanged as \texttt{sources/153020/source0.tex}.
\begin{lemma}\label{lem:scorelhs}
    Define $\mathcal{F}:=\{F\in \mathbb{R}^{k\times k}:\|F\|\leq 1\}$. Let $\left\{\bx^{(i)}\right\}_{i \in [n]}$ be i.i.d samples from a subgaussian$\bb{\sigma}$ distribution. We have: %Then, for $\forall \bt \in \mathbb{R}^{d}, \norm{\bt} = 1$, we have
    \bas{
        \sup_{F\in \mathcal{F}}\left|\frac{1}{n}\sum_{j=1}^{n}\exp(i\bt^TF\bx^{\bb{j}})-\E_{\bx}\bbb{\exp(i \bt^T F\bx)}\right| = O_P\bb{\nt\sqrt{\frac{k\Tr(S)\log\bb{nC_{t}}}{n}}}
    }
    where $C_{t} := \max(1,
    \nts\Tr(S))$.
\end{lemma}

\begin{lemma}\label{lem:scorerhs}
Let $\mathcal{F}=\{F\in \mathbb{R}^{k\times k}:\|F\|\leq 1\}$. We have:
\bas{
    \sup_{F\in \mathcal{F}}\left|\psi(\bt,F|\hat{P})-\psi(\bt,F|P)\right|\leq  O_P\bb{\nt\sqrt{\frac{k^2\|S\|\log n}{n}}}
}
\end{lemma}

\begin{proof}
%Let $F_j$ denote the $j^{th}$ row of $F$. $\mathcal{F}_j:=\{\bu\in \mathbb{R}^k:\|\bu\|\leq 1\}$. 
Let $\mathcal{B}_{k} := \left\{\bu \in \mathbb{R}^{k}, \norm{\bu}\leq1\right\}$
Now define $\psi_j(\bt,F;\bx)=\E[\exp(it_j F_j^T\bx)]$. Also define $f^{(j)}_X(F)=\frac{1}{n}\sum_i\epsilon_i \psi_j(\bt,F;\bx^{\bb{i}})$. 
% The only things that change are the distance, diameter, and covering number. Define 
% \bas{
%     d(F_j,F_j')^2:=\sum_i (\psi_j(\bt,F;\bx^{(i)})-\psi_j(\bt,F';\bx^{(i)}))^2
% }
%Also define $\mathcal{F}:=\{F:\|F\|_F\leq 1\}$.
% First, note that $\psi_j(\bt,F;\bx)$ is Lipschitz over $F$. Let $\be_j$ denote a vector with all zeros except one at the $j^{th}$ position.
% \bas{
%     \|\nabla \psi_j(\bt,F;\bx)\|_F = \|t_j \be_j \bx^T\|_F|-\sin(t_j (F\bx)_j)+i\cos(t_j (F\bx)_j)|\leq \|\be_{j}\bx^{T}\|_{F}|t_j|
% }
% Thus,
% \bas{
% D:=\max_{F,F'}d(F_j,F_j')\leq \sqrt{n}{\max_i\|\bx^{(i)}\|}|t_j|
% }
% Next, observe that, with probability $1-\delta$,
% \bas{
% &n\E\sup_{F\in\mathcal{F}^{(j)}}|f^{(j)}_X(F)-f^{(j)}_X(F')|\\
% %&\leq 2\E \sup_{\|F\|\leq 1}|f_X(F)-f_X(F')|+\exp(-c_1 n\delta^2)\\
% &\leq 2\E\sup_{d(F_j,F'_j)\leq \epsilon}|f^{(j)}_X(F)-f^{(j)}_X(F')|+2\E[D]\sqrt{\log N(\epsilon;\mathcal{F}_j,d_F)}\\
% &\leq |t_j|\E\max_i\|\bx^{(i)}\|\bb{\epsilon\sqrt{n} + 4c_2 \sqrt{nk\log (1/\epsilon)}}\\
% &= O\bb{\sqrt{nk^2\log n}}
% }
% The last step follows because $\bx^{(i)}\sim \text{subgaussian}(\sigma)$.
Thus, using the same argument as in Lemma~\ref{lem:scorelhs} for $\bt \leftarrow \nt\be_{j}$ and $\bx^{\bb{i}} \leftarrow t_{j}\bx^{\bb{i}}$, 
\bas{\sup_{F\in \mathcal{F}}|\psi_j(\bt,F|\hat{P})-\psi_j(\bt,F|P)|=O_P\bb{|t_{j}|\sqrt{\frac{k\|S\|\log n}{n}}}
}
Finally, we see that:
% \rd Syamantak, please check: \bk
\bas{
\sup_{F\in \mathcal{F}}\left|\psi(\bt,F|\hat{P})-\psi(\bt,F|P)\right| &\leq \sum_j  O_P\bb{|t_j|\sqrt{\frac{k\|S\|\log n}{n}}}\\
&= O_P\bb{\sqrt{\frac{k\|S\|\log n}{n}}}\bb{\sum_{j=1}^{k}|t_j|} 
&= \nt O_P\bb{\sqrt{\frac{k^2\|S\|\log n}{n}}}
}
The above is true because, for $|a_i|,|b_i|\leq 1$, $i=1,\dots, k$,
\bas{
\left|\prod_{i=1}^k a_i-\prod_{i=1}^k b_i\right|=\left|\sum_{j=0}^{k-1} \prod_{i\leq j}
b_i (a_{j+1}-b_{j+1})\prod_{i=j+2}^k
a_i\right| \leq \sum_{j=0}^{k-1}|a_{j+1}-b_{j+1}| %-\prod_{i=2}^k a_i+b_1(a_2-b_2)\prod_{i=3}^k a_i+\dots+b_1\dots (a_k-b_k)
}
\end{proof}

\end{document}
